
\documentclass[10pt]{article}
\usepackage[utf8]{inputenc}
\usepackage[T1]{fontenc}
\usepackage{amsmath, amssymb, xcolor, geometry, enumitem, multicol, tcolorbox}

\definecolor{mainblue}{RGB}{0, 102, 204}
\geometry{a4paper, margin=1.2cm, top=1cm, bottom=3cm, footskip=1.2cm}

\begin{document}
\enlargethispage{2.5cm}

% Modern Header
\begin{tcolorbox}[colback=mainblue!5, colframe=mainblue, sharp corners, boxrule=0.5pt]
    \small \textbf{Unit:} undefined \hfill \textbf{Name:} \rule{3cm}{0.4pt} \\
    \small \textbf{Topic:} Variables and Expressions / Order of Operations \hfill \textbf{Date:} \rule{2cm}{0.4pt} \\
    \centering \textbf{\large Foundation (Level -1)}
\end{tcolorbox}

\vspace{0.2cm}

% MCQ Section
\begin{multicols}{2}
\begin{enumerate}[label=\textbf{Q\arabic*.}, leftmargin=*, itemsep=6pt, parsep=0pt]

  \item \small What is the variable in the expression $2x$?
  \begin{enumerate}[label=(\Alph*), leftmargin=0.6cm, nosep, itemsep=2pt]
    \item 2
    \item x
    \item 2x
    \item +
    \item None of the above
  \end{enumerate}

  \item \small Identify the coefficient in the term $4y$?
  \begin{enumerate}[label=(\Alph*), leftmargin=0.6cm, nosep, itemsep=2pt]
    \item y
    \item 4
    \item 4y
    \item None of the above
    \item 0
  \end{enumerate}

  \item \small What is the constant term in the expression $3x + 2$?
  \begin{enumerate}[label=(\Alph*), leftmargin=0.6cm, nosep, itemsep=2pt]
    \item 3x
    \item 2
    \item 3
    \item +
    \item None of the above
  \end{enumerate}

  \item \small Identify the operator in $7 - 3$?
  \begin{enumerate}[label=(\Alph*), leftmargin=0.6cm, nosep, itemsep=2pt]
    \item 7
    \item 3
    \item -
    \item None of the above
    \item 0
  \end{enumerate}

  \item \small What is the variable in the expression $z + 1$?
  \begin{enumerate}[label=(\Alph*), leftmargin=0.6cm, nosep, itemsep=2pt]
    \item 1
    \item z
    \item z + 1
    \item +
    \item None of the above
  \end{enumerate}

  \item \small Identify the coefficient in the term $x$?
  \begin{enumerate}[label=(\Alph*), leftmargin=0.6cm, nosep, itemsep=2pt]
    \item 1
    \item x
    \item None of the above
    \item 0
    \item -1
  \end{enumerate}

  \item \small What operation is represented by the symbol $+$?
  \begin{enumerate}[label=(\Alph*), leftmargin=0.6cm, nosep, itemsep=2pt]
    \item Subtraction
    \item Addition
    \item Multiplication
    \item Division
    \item None of the above
  \end{enumerate}

  \item \small Identify the constant term in $2x - 5$?
  \begin{enumerate}[label=(\Alph*), leftmargin=0.6cm, nosep, itemsep=2pt]
    \item $2x$
    \item -5
    \item 2
    \item None of the above
    \item 5
  \end{enumerate}

\end{enumerate}
\end{multicols}

\vspace{-0.2cm}
\hrule
\vspace{0.2cm}
\textbf{\small Open Ended Questions (FRQ)}
\begin{enumerate}[label=\textbf{Q\arabic*.}, start=9, itemsep=5pt, topsep=0pt]

  \item \small Draw a diagram to represent the expression $3x + 2$. Label the variable and constant terms. \vspace{2cm}

  \item \small Explain the concept of a variable in an expression, using $x$ as an example. Provide an illustration. \vspace{2cm}

\end{enumerate}

\vfill

% Chef's Tips Footer
\begin{tcolorbox}[colback=blue!5, colframe=mainblue!50, title=\small \textbf{Chef's Tips}, fonttitle=\bfseries, bottom=1mm]
\begin{itemize}[noitemsep, topsep=2pt, leftmargin=0.5cm]
  \item \scriptsize Focus on basic conceptual definitions and single-step identification for Level -1 students.\item \scriptsize Ensure students understand the difference between variables, coefficients, and constants in expressions.\item \scriptsize Use visual aids and diagrams to help students comprehend abstract concepts, such as variables and expressions.
\end{itemize}
\end{tcolorbox}

\end{document}
  